\[ %%%%%%%%%%%%%%%%%%%%%%%%%%%%%%%%%%%%%%%%%%%%%%%%%%%%%%%%%%%%%%%%%%%%%%%%%%%%%%%% \subsection{Physics-Guided Graph Evolution Equation} \label{sec:graph_evolution_equation} %%%%%%%%%%%%%%%%%%%%%%%%%%%%%%%%%%%%%%%%%%%%%%%%%%%%%%%%%%%%%%%%%%%%%%%%%%%%%%%% The central mathematical formulation of PHYSGEN is a physics-guided graph evolution equation describing the temporal evolution of a physical system as a dynamic interaction graph. Unlike conventional neural forecasting approaches that directly learn mappings between consecutive observations, PHYSGEN learns an evolution operator acting on the complete physical graph state. The general graph evolution process is formulated as: \begin{equation} \boxed{ G_{t+1} = \Phi_{\theta} ( G_t, P_t, C_t ) } \end{equation} where: \begin{itemize} \item $G_t=(V_t,E_t,X_t,P_t)$ is the dynamic physical graph; \item $P_t$ represents physical parameters; \item $C_t$ represents physical constraints; \item $\Phi_{\theta}$ is the physics-guided evolution operator. \end{itemize} %%%%%%%%%%%%%%%%%%%%%%%%%%%%%%%%%%%%%%%%%%%%%%%%%%%%%%%%%%%%%%%%%%%%%%%%%%%%%%%% \subsubsection{Decomposition of the Evolution Operator} %%%%%%%%%%%%%%%%%%%%%%%%%%%%%%%%%%%%%%%%%%%%%%%%%%%%%%%%%%%%%%%%%%%%%%%%%%%%%%%% The PHYSGEN evolution operator is decomposed into four coupled components: \begin{equation} \Phi_{\theta} = \Phi_{stab} \circ \Phi_{dyn} \circ \Phi_{phys} \circ \Phi_{enc}, \end{equation} where: \begin{itemize} \item $\Phi_{enc}$ performs physical graph encoding; \item $\Phi_{phys}$ performs physics-guided interaction propagation; \item $\Phi_{dyn}$ models latent physical evolution; \item $\Phi_{stab}$ applies adaptive stability regulation. \end{itemize} The complete state transition can therefore be expressed as: \begin{equation} G_t \rightarrow H_t \rightarrow H_{t+1} \rightarrow \hat{G}_{t+1}. \end{equation} %%%%%%%%%%%%%%%%%%%%%%%%%%%%%%%%%%%%%%%%%%%%%%%%%%%%%%%%%%%%%%%%%%%%%%%%%%%%%%%% \subsubsection{Physics-Guided Latent Graph Evolution} %%%%%%%%%%%%%%%%%%%%%%%%%%%%%%%%%%%%%%%%%%%%%%%%%%%%%%%%%%%%%%%%%%%%%%%%%%%%%%%% Let the latent graph representation be: \begin{equation} Z_t = Enc_{\theta}(G_t). \end{equation} The latent evolution is defined as: \begin{equation} Z_{t+1} = Z_t + \Delta t [ F_{\theta}(Z_t,G_t) + F_{phys}(Z_t,P_t) ]. \end{equation} The first term represents the learned nonlinear dynamics: \begin{equation} F_{\theta}(Z_t,G_t), \end{equation} while the second term introduces physics-guided corrections: \begin{equation} F_{phys}(Z_t,P_t). \end{equation} This formulation allows PHYSGEN to combine expressive neural dynamics with physical consistency. %%%%%%%%%%%%%%%%%%%%%%%%%%%%%%%%%%%%%%%%%%%%%%%%%%%%%%%%%%%%%%%%%%%%%%%%%%%%%%%% \subsubsection{Graph Topology Evolution} %%%%%%%%%%%%%%%%%%%%%%%%%%%%%%%%%%%%%%%%%%%%%%%%%%%%%%%%%%%%%%%%%%%%%%%%%%%%%%%% In many physical systems, interaction topology is not static. Therefore, PHYSGEN models graph connectivity evolution: \begin{equation} E_{t+1} = \Gamma_{\theta} ( E_t, X_t, P_t ). \end{equation} The complete graph transition becomes: \begin{equation} G_{t+1} = ( V_{t+1}, E_{t+1}, X_{t+1}, P_{t+1} ). \end{equation} This enables modeling of systems with: \begin{itemize} \item changing spatial relationships; \item adaptive interaction networks; \item structural transitions; \item evolving physical configurations. \end{itemize} %%%%%%%%%%%%%%%%%%%%%%%%%%%%%%%%%%%%%%%%%%%%%%%%%%%%%%%%%%%%%%%%%%%%%%%%%%%%%%%% \subsubsection{Physics Constraint Projection} %%%%%%%%%%%%%%%%%%%%%%%%%%%%%%%%%%%%%%%%%%%%%%%%%%%%%%%%%%%%%%%%%%%%%%%%%%%%%%%% After latent evolution, the predicted state is projected onto the physically admissible manifold. The projection operation is defined as: \begin{equation} \tilde{Z}_{t+1} = \Pi_{\mathcal{M}_{phys}} ( Z_{t+1} ), \end{equation} where: \begin{equation} \Pi_{\mathcal{M}_{phys}} \end{equation} represents a physics-consistency projection operator. The final predicted state becomes: \begin{equation} \hat{X}_{t+1} = Dec_{\theta} ( \tilde{Z}_{t+1} ). \end{equation} %%%%%%%%%%%%%%%%%%%%%%%%%%%%%%%%%%%%%%%%%%%%%%%%%%%%%%%%%%%%%%%%%%%%%%%%%%%%%%%% \subsubsection{Unified PHYSGEN Evolution Equation} %%%%%%%%%%%%%%%%%%%%%%%%%%%%%%%%%%%%%%%%%%%%%%%%%%%%%%%%%%%%%%%%%%%%%%%%%%%%%%%% Combining all components, the complete PHYSGEN transition equation is: \begin{equation} \boxed{ X_{t+1} = Dec_{\theta} \left[ \Pi_{\mathcal{M}_{phys}} \left( Z_t + \Delta t ( F_{\theta}(Z_t,G_t) + F_{phys}(Z_t,P_t) ) + C_{stab} \right) \right] } \end{equation} This equation represents the complete physics-guided evolution mechanism of the framework. %%%%%%%%%%%%%%%%%%%%%%%%%%%%%%%%%%%%%%%%%%%%%%%%%%%%%%%%%%%%%%%%%%%%%%%%%%%%%%%% \subsubsection{Continuous-Time Interpretation} %%%%%%%%%%%%%%%%%%%%%%%%%%%%%%%%%%%%%%%%%%%%%%%%%%%%%%%%%%%%%%%%%%%%%%%%%%%%%%%% The discrete evolution equation can be interpreted as a neural approximation of a continuous physical operator: \begin{equation} \frac{dZ(t)}{dt} = F_{\theta}(Z(t),G(t)) + F_{phys}(Z(t),P). \end{equation} Thus, PHYSGEN can be viewed as a physics-guided neural dynamical system operating on adaptive graph representations. %%%%%%%%%%%%%%%%%%%%%%%%%%%%%%%%%%%%%%%%%%%%%%%%%%%%%%%%%%%%%%%%%%%%%%%%%%%%%%%% \subsubsection{Relation to Scientific Machine Learning} %%%%%%%%%%%%%%%%%%%%%%%%%%%%%%%%%%%%%%%%%%%%%%%%%%%%%%%%%%%%%%%%%%%%%%%%%%%%%%%% The proposed formulation establishes a connection between several scientific machine learning paradigms: \begin{itemize} \item Graph neural simulators provide relational representation; \item Neural operators provide function-space evolution learning; \item Physics-informed methods provide constraint mechanisms; \item Neural dynamical systems provide temporal evolution modeling. \end{itemize} PHYSGEN integrates these concepts into a unified graph-based physics-guided evolution framework. %%%%%%%%%%%%%%%%%%%%%%%%%%%%%%%%%%%%%%%%%%%%%%%%%%%%%%%%%%%%%%%%%%%%%%%%%%%%%%%% \subsubsection{Summary} %%%%%%%%%%%%%%%%%%%%%%%%%%%%%%%%%%%%%%%%%%%%%%%%%%%%%%%%%%%%%%%%%%%%%%%%%%%%%%%% The physics-guided graph evolution equation defines the mathematical core of PHYSGEN. By learning a constrained evolution operator over dynamic physical graphs, the framework moves beyond static prediction toward adaptive modeling of nonlinear physical processes with improved long-horizon stability and physical consistency. \] 
